% STATUS: DRAFT
\chapter{KMS equilibrium}\label{ch:kms}

\begin{physmotiv}
The Gibbs state $\rho=\ee^{-\beta H}/Z$ is the foundation of statistical mechanics, but in infinite volume neither $\ee^{-\beta H}$ nor $Z=\Tr\ee^{-\beta H}$ makes sense (\cref{ch:quasilocal}). What survives is a \emph{property} of the Gibbs state that can be stated purely in terms of expectation values and the dynamics: the Kubo--Martin--Schwinger (KMS) condition. Haag, Hugenholtz and Winnink proposed taking it as the \emph{definition} of equilibrium. The same condition, applied to an arbitrary state rather than a Gibbs state, is what Tomita--Takesaki theory (\cref{ch:tomita}) turns into a canonical dynamics: every faithful state is an equilibrium state of its own modular flow.
\end{physmotiv}

\section{The finite-dimensional origin}

Let $\A=\BH$ with $\dim\HH<\infty$, $\alpha_t(a)=\ee^{\ii tH}a\,\ee^{-\ii tH}$ and $\omega(a)=\Tr(\ee^{-\beta H}a)/Z$. Extend $\alpha_t$ to complex time: $\alpha_{z}(a)=\ee^{\ii zH}a\,\ee^{-\ii zH}$, entire in $z$ for a finite matrix. Then, using only cyclicity of the trace,
\begin{equation}
\omega\big(a\,\alpha_{\ii\beta}(b)\big)=\tfrac1Z\Tr\big(\ee^{-\beta H}a\,\ee^{-\beta H}b\,\ee^{\beta H}\big)=\tfrac1Z\Tr\big(a\,\ee^{-\beta H}b\big)=\omega(ba).\label{eq:kmsfinite}
\end{equation}
This is the \emph{KMS condition}. It encodes the cyclic property of the trace in the presence of the Boltzmann weight and is equivalent to the statement that correlation functions are periodic in imaginary time with period $\beta$: with $F(t)=\omega(a\,\alpha_t(b))$ one has $F(t+\ii\beta)=\omega(\alpha_t(b)\,a)$.

\section{The KMS condition in general}

\begin{definition}
Let $\alpha_t$ be a one-parameter group of automorphisms of a \cstar-algebra (or von Neumann algebra) $\A$. A state $\omega$ is a \emph{$(\alpha,\beta)$-KMS state}, $\beta\in\R\setminus\{0\}$, if for every $a,b\in\A$ there is a function $F_{a,b}$, continuous and bounded on the closed strip $\{0\le\mathrm{Im}\,z\le\beta\}$ (for $\beta<0$: $\beta\le\mathrm{Im}\,z\le0$) and analytic in its interior, with
\begin{equation}
F_{a,b}(t)=\omega\big(a\,\alpha_t(b)\big),\qquad F_{a,b}(t+\ii\beta)=\omega\big(\alpha_t(b)\,a\big).
\end{equation}
Equivalently, $\omega\big(a\,\alpha_{\ii\beta}(b)\big)=\omega(ba)$ for $b$ in a dense set of $\alpha$-analytic elements.
\end{definition}

The proposal (Haag--Hugenholtz--Winnink): \emph{an infinite system in thermal equilibrium at inverse temperature $\beta$ is described by an $(\alpha,\beta)$-KMS state}. This works because
\begin{itemize}
\item Gibbs states of finite systems satisfy it exactly, and weak-$^*$ limits of finite-volume Gibbs states (\cref{ch:quasilocal}) satisfy it in the limit;
\item KMS states are automatically $\alpha$-invariant, $\omega\circ\alpha_t=\omega$ (stationary);
\item for fixed $\beta$ the KMS states form a simplex; the extremal points are the pure thermodynamic phases, and they are factor states (\cref{ch:quasilocal}).
\end{itemize}
The limit $\beta\to\infty$ gives ground states; $\beta=0$ gives tracial states.

\begin{whatwrong}
Sign conventions for $\beta$ differ in the literature (the sign of $t$ in $\alpha_t$, and whether the condition is written $\omega(a\alpha_{\ii\beta}(b))=\omega(ba)$ or $\omega(\alpha_{-\ii\beta}(a)b)=\omega(ba)$). We use $\alpha_t(a)=\ee^{\ii tH}a\,\ee^{-\ii tH}$ and the condition $\omega(a\alpha_{\ii\beta}(b))=\omega(ba)$, which gives Gibbs states $\ee^{-\beta H}$ with $\beta>0$ for $H$ bounded below. Every statement of the form ``the Rindler vacuum is KMS at $\beta=2\pi$'' needs this fixed before it means anything.
\end{whatwrong}

\section{Examples}

\subsection{Free Bose gas}
For the Weyl algebra of the free field with one-particle Hamiltonian $\omega$ (an operator, as in \cref{ch:gns}), the quasi-free state
\begin{equation}
\omega_\beta\big(W(f)\big)=\exp\Big[-\tfrac12\big\langle Kf,\coth\!\big(\tfrac{\beta\omega}{2}\big)Kf\big\rangle\Big]
\end{equation}
is the unique $\beta$-KMS state for the free dynamics. The factor $\coth(\beta\omega/2)=1+2n_B(\omega)$ is Bose--Einstein. As $\beta\to\infty$ it tends to $1$ and one recovers the vacuum state. In infinite volume the representation is inequivalent to the vacuum (\cref{ch:inequiv}) and the GNS space is $\mathcal F\otimes\mathcal F$ (the thermofield double).

\subsection{Spin chains}
For non-interacting spins $H=\frac h2\sum\sigma^z_j$ the KMS state is the product state of \cref{ch:quasilocal}, $\omega_\beta(\sigma^z_j)=-\tanh(\beta h/2)$. For short-range interacting chains in one dimension there is a unique KMS state at each $\beta$ (no phase transition at $T>0$, Araki), while in higher dimensions there may be several (spontaneous magnetization).

\subsection{The Unruh/Rindler state}
For a relativistic QFT, the vacuum $\Omega$ restricted to the algebra $\A(R)$ of the right wedge is a KMS state for the \emph{Lorentz boosts} $\alpha_s=\Ad\,U(\Lambda(s))$ with inverse temperature $\beta=2\pi$ in rapidity units (for boosts $\Lambda(s)$ of rapidity $s$ that move the wedge forward in time, so that $\alpha_s=\Ad\ee^{\ii sB}$ with $B$ the boost generator); this is the Bisognano--Wichmann theorem (\cref{ch:rs_bw}). Converting from rapidity to proper time for an observer at acceleration $a$ gives the Unruh temperature $T=a/2\pi$. The point of the algebraic formulation: no wedge Hamiltonian or $\Tr$ is used, only the KMS condition. The same structure for the static patch of de Sitter space will give the Gibbons--Hawking temperature (\cref{ch:static_patch}).

\section{Stability, passivity, and the physics of KMS}

Why should KMS be the ``right'' notion of equilibrium? Two physical criteria both lead to it.

\textbf{Passivity (Pusz--Woronowicz).} Imagine perturbing the system by a time-dependent Hamiltonian which is switched on and then off again, so that the net effect is a unitary $U$ in the algebra. A state is \emph{passive} if no energy can be extracted by such a cyclic process. The state is \emph{completely passive} if this remains true for any number of independent copies of the system. A state is completely passive iff it is a ground state or a KMS state at some $\beta\ge0$. This is the second law (Kelvin's formulation) for infinite systems. It also shows why states with $\beta<0$ (population inversions) are \emph{not} stable.

\textbf{Stability under perturbations (Araki).} If $\omega$ is a KMS state and $\alpha^P$ is the dynamics perturbed by a local self-adjoint $P\in\A$, there is a unique perturbed KMS state $\omega^P$ at the same temperature, constructed by a norm-convergent expansion in the analytic continuation (the algebraic version of $\ee^{-\beta(H+P)}$). So KMS states are stable under local perturbations, and the response (Kubo formula) can be derived entirely algebraically.

\section{The punchline for modular theory}

Suppose a faithful normal state $\omega$ on a von Neumann algebra $\M$ is given and no dynamics is specified. One can ask: is there a one-parameter group $\sigma_t$ of $\M$ for which $\omega$ is KMS? In finite dimensions, $\omega=\Tr(\rho\,\cdot)$, take $\sigma_t(a)=\rho^{it}a\rho^{-it}$. Then $\sigma_t=\alpha_{-\beta t}$ with $\rho=\ee^{-\beta H}/Z$, and the KMS condition becomes (using $\beta=-1$ in the sense above)
\begin{equation}
\omega\big(a\,\sigma_{-\ii}(b)\big)=\omega(ba).\label{eq:kmsmod}
\end{equation}
\begin{theorem}[Takesaki, previewed]
For any von Neumann algebra $\M$ and faithful normal state $\omega$ there is a \emph{unique} strongly continuous one-parameter group of automorphisms $\sigma^\omega_t$ of $\M$ for which $\omega$ satisfies \eqref{eq:kmsmod}, the \emph{modular group}. It is trivial iff $\omega$ is a trace.
\end{theorem}
Think of it as: \emph{every faithful state is a thermal state with respect to some canonical Hamiltonian, at unit negative temperature $\beta=-1$}, the minus sign being a convention of the time direction. For a Gibbs state with Hamiltonian $H$ the modular Hamiltonian is $\beta H$, up to a constant. The modular flow is thus ``physical time'' multiplied by $-\beta$, and the information about temperature has been absorbed into the time unit. We prove it in \cref{ch:tomita}.

\begin{keyidea}
KMS replaces ``$\ee^{-\beta H}$'' by a condition on correlation functions: analytic continuation in time by $\ii\beta$ turns $\omega(a\,\alpha_t(b))$ into $\omega(\alpha_t(b)\,a)$. It defines equilibrium in infinite volume, is equivalent to passivity, and applied to an arbitrary faithful state it \emph{defines} the modular flow.
\end{keyidea}

\section*{Exercises}
\begin{exercise}
Verify \eqref{eq:kmsfinite}. Show that for the Gibbs state $F(t)=\omega(a\alpha_t(b))$ extends to an entire function with $F(t+\ii\beta)=\omega(\alpha_t(b)a)$. Check that $\omega\circ\alpha_t=\omega$ (stationarity) follows from KMS.
\end{exercise}
\begin{exercise}
For a two-level system with $H=\frac h2\sigma^z$ and $b=\sigma^+$, $a=\sigma^-$, compute $\omega(a\,\alpha_t(b))$ and $\omega(\alpha_t(b)\,a)$ in the Gibbs state, and verify the strip condition explicitly: find the analytic function and check $F(t+\ii\beta)$. Determine the ratio $\omega(\sigma^-\sigma^+)/\omega(\sigma^+\sigma^-)$ (detailed balance).
\end{exercise}
\begin{exercise}
Show that if $\omega$ is a $(\alpha,\beta)$-KMS state with $\beta\ne0$ and $\omega$ is a vector state of a vector $\Omega$ separating for $\A''$, then $t\mapsto\alpha_t$ is implemented in the GNS representation by $\Delta_\Omega^{-it/\beta}$ for a positive operator $\Delta_\Omega$ (anticipating \cref{ch:tomita}). In the finite-dimensional case, find $\Delta$ on $\HH\otimes\overline\HH$.
\end{exercise}
\begin{exercise}
Show that $\beta\to0$ gives a tracial state: the KMS condition at $\beta=0$ reads $\omega(ab)=\omega(ba)$. Show that the infinite-temperature state of the spin chain is KMS at $\beta=0$ for every dynamics.
\end{exercise}
